\section{Easy, Medium, Hard: por qué una sola decodificación subestima la capacidad de la distribución}

El ponente usa tres HumanEval problem para mostrar las diferencias de orden de magnitud en pass@1. El valor de las slides no está en memorizar la implementación concreta, sino en observar cómo la prompt ambiguity, la algorithm depth y la rare correctness modifican la sampling strategy.

\subsection{Easy problem}

Esta sección examina primero el easy case para establecer la línea base en la que «las muestras correctas se ubican en la región de alta probabilidad», y después aumenta gradualmente la dificultad de algorithm y de boundary. Al leer el código, primero se enumera a mano la specification a partir del docstring y luego se observa si el modelo cubre los casos normal y los edge cases solo con una plantilla común.

\lecturefigure{slide-34-easy-problem.jpg}{La implementación correcta del Easy problem se ubica en la región de alta probabilidad de la distribución del modelo.}{Video oficial 38:03.}

\begin{knowledgebox}{Lectura de la figura: un pass@1 alto también exige leer la specification}
Cuando el código es corto, el pattern es común y hay pocos boundary, el greedy o un sample a baja temperatura aciertan con facilidad. Primero se verifica que el prompt esté completo y después que el candidate maneje empty/negative/duplicate. Una puntuación alta puede provenir de plantillas frecuentes en el entrenamiento y no implica que el modelo comprenda todas las variantes.
\end{knowledgebox}

\subsection{Medium problem}

A continuación se observa el medium case: la estructura principal puede ser correcta, pero basta un branch o un off-by-one para que fallen todas las hidden tests. Este tipo de problema es el que mejor muestra que la fluidez en lenguaje natural está desvinculada de la corrección en ejecución, y también el más adecuado para analizar si múltiples sampling producen errores complementarios.

\lecturefigure{slide-35-medium-problem.jpg}{El Medium problem exige más pasos o más manejo de casos límite, y la probabilidad de obtener una muestra correcta baja notablemente.}{Video oficial 38:36.}

\begin{knowledgebox}{Lectura de la figura: el error puede concentrarse en un solo branch}
Los candidatos suelen tener ya el esqueleto correcto, pero fallan en el rango del ciclo, en el orden de las condiciones o en el state update. En ese caso, múltiples sampling pueden explorar la variante correcta; también se puede mejorar la selección con test feedback, static analysis o un repair loop. El benchmark solo cuenta el pass final y no explica el tipo de error, por lo que se necesita una taxonomy adicional.
\end{knowledgebox}

\subsection{Hard problem}

Al pasar al hard case, la pregunta ya no es «si la salida por defecto es correcta», sino si el program correcto tiene en la distribución suficiente probability mass para que la search lo encuentre. Si la mayoría de los samples comparten el mismo algorithm erróneo, aumentar $k$ solo repite el fracaso; solo si los errores son diversos puede el verifier filtrar un candidato correcto de la cola larga.

\lecturefigure{slide-36-hard-problem.jpg}{La tasa de acierto con una sola muestra en el Hard problem es muy baja, pero la distribución todavía puede contener implementaciones correctas.}{Video oficial 39:00.}

\begin{knowledgebox}{Lectura de la figura: un pass@1 bajo no equivale a capacidad nula}
Si pass@1 es del orden de milésimas, una demo de un solo intento casi siempre falla; una gran cantidad de samples todavía puede encontrar un programa correcto. El sistema debe decidir si puede asumir el costo de generar, ejecutar y filtrar. Sin un verifier confiable, large-k solo produce más código que no puede auditarse.
\end{knowledgebox}

\subsection{Por qué se necesita una sampling-aware metric}

Ahora se pasa de los ejemplos de un solo problema a la distribución de todo el benchmark. El histograma muestra que algunos problemas se resuelven casi siempre y que otros presentan una cola larga; el pass@1 promedio oculta la diferencia entre «basta con pocos candidatos para resolverlo» y «no se resuelve ni con muchas más muestras». La pregunta que conviene llevar es: a qué problemas debe asignarse el generation budget adicional y cuándo conviene mejorar el modelo en lugar de seguir muestreando.

\lecturefigure{slide-37-need-pass-at-k-hist.jpg}{La probabilidad de éxito por problema tiene una cola larga; una sola puntuación promedio oculta las diferencias en lo que puede obtenerse mediante muestreo.}{Video oficial 40:30.}

\begin{knowledgebox}{Lectura de la figura: eje horizontal, tasa de éxito por problema; eje vertical, número de problemas}
Si muchos problemas se concentran cerca de 0 o de 1, el desempeño del modelo es bimodal; la región intermedia es la que más puede mejorarse mediante sampling. Hay que reportar la per-problem distribution y no solo la mean. También conviene revisar cómo se desplaza la mass con distintos model size y si los problemas difíciles obtienen realmente una probabilidad distinta de cero.
\end{knowledgebox}

\lecturefigure{slide-38-need-pass-at-k-scatter.jpg}{La mejora de pass@1 a pass@k no es uniforme entre problemas.}{Video oficial 40:57.}

\begin{knowledgebox}{Lectura de la figura: que un punto se aleje de la diagonal indica sampling value}
Si un punto sube notablemente en el eje de k alto, el muestreo diverso puede encontrar rare correct programs; los puntos que siguen pegados a cero indican que la distribution casi no tiene masa correcta. El sampling no lo resuelve todo: solo amplifica la probability mass existente y no puede crear de la nada un algorithm que el modelo nunca aprendió.
\end{knowledgebox}

\subsection{Resumen de la sección}

La dificultad del problema cambia la posición de las muestras correctas en la distribución. Pass@1 mide el acierto por defecto en un solo intento, y pass@k mide el potencial de búsqueda de un sistema de generación y filtrado; ambos requieren una interpretación per-problem.
